\documentclass[10pt]{article}
\usepackage[T1]{fontenc}
\usepackage[utf8]{inputenc}
\usepackage[french]{babel}
\usepackage{amsmath,amssymb,mathrsfs,mathpazo}
\usepackage{geometry}
\geometry{a4paper,margin=2cm}
\usepackage{graphicx,array,tabularx,tabularray,colortbl,multirow,multicol,booktabs,enumitem,pifont,ulem,fancybox,diagbox,cancel}
\usepackage{pstricks,pstricks-add,pst-plot,pst-tree,pst-node,pst-text,pst-func,pst-math,pst-bspline,pst-3dplot,pst-eucl,pst-solides3d}
\let\pscalculate\undefined
\usepackage{pst-bezier}
\usepackage{tikz,pgfplots,tkz-tab}
\usetikzlibrary{arrows,arrows.meta,patterns,calc,positioning,trees,decorations.markings,shapes.geometric,shapes.misc,shapes.arrows,intersections,decorations.pathmorphing,decorations.pathreplacing,backgrounds,shadows,babel}
\pgfplotsset{compat=1.18}
\usepgfplotslibrary{fillbetween}
\usepackage[np]{numprint}
\usepackage{listings,makecell,pas-tableur,twoopt}
\let\barre\undefined
\usepackage[pstricks]{tablvar}
\definecolor{darkgreen}{RGB}{0,100,0}
\providecommand{\red}{\color{red}}
\providecommand{\blue}{\color{blue}}
\providecommand{\green}{\color{green}}
\providecommand{\black}{\color{black}}
\providecommand{\white}{\color{white}}
\providecommand{\gray}{\color{gray}}
\providecommand{\cyan}{\color{cyan}}
\providecommand{\magenta}{\color{magenta}}
\providecommand{\yellow}{\color{yellow}}
\providecommand{\violet}{\color{violet}}
\providecommand{\orange}{\color{orange}}
\providecommand{\euro}{\texteuro}
\providecommand{\EUR}{\texteuro}
\providecommand{\ieme}{\textsuperscript{e}}
\providecommand{\degres}{\ensuremath{^\circ}}
\providecommand{\degree}{\ensuremath{^\circ}}
\providecommand{\vv}[1]{\overrightarrow{#1}}
\providecommand{\Coord}[1]{#1}
\providecommand{\decosix}{}
\providecommand{\starredbullet}{\ensuremath{\bullet}}
\pagestyle{empty}
\setlength{\parindent}{0pt}
\usepackage{adjustbox}\DeclareUnicodeCharacter{2236}{:}\DeclareUnicodeCharacter{2190}{\ensuremath{\leftarrow}}
\providecommand{\R}{\mathbb{R}}
\providecommand{\N}{\mathbb{N}}
\providecommand{\Z}{\mathbb{Z}}
\providecommand{\Q}{\mathbb{Q}}
\providecommand{\D}{\mathbb{D}}
\providecommand{\C}{\mathbb{C}}
\providecommand{\e}{{\mathrm{e}}}
\providecommand{\vect}[1]{\overrightarrow{#1}}
\providecommand{\vectt}[1]{\overrightarrow{\mathrm{#1}}}
\providecommand{\vvt}[1]{\overrightarrow{\mathrm{#1}}}
\providecommand{\np}[2][]{{#2~#1}}
\providecommand{\vv}[1]{{\overrightarrow{#1}}}
\renewcommand{\d}{{\mathrm{d}}}
\providecommand{\degres}{{^\circ}}
\providecommand{\degre}{{^\circ}}
\providecommand{\blue}{\color{blue}}
\providecommand{\red}{\color{red}}
\providecommand{\enskip}{\quad}
\providecommand{\llbracket}{\lbrack\!\lbrack}
\providecommand{\rrbracket}{\rbrack\!\rbrack}
\providecommand{\euro}{{\text{€}}}
\providecommand{\no}{{\text{n°}}}
\providecommand{\interval}[2]{{[#1;#2]}}
\providecommand{\intervFO}[2]{{[#1;#2[}}
\providecommand{\dotfill}{\ldots}

\providecommand{\ds}{\displaystyle}
\providecommand{\textding}[1]{#1}
\providecommand{\Oij}{\(\mathrm{O};\vec{\imath},\vec{\jmath}\)}
\providecommand{\Oijk}{\(\mathrm{O};\vec{\imath},\vec{\jmath},\vec{k}\)}
\providecommand{\Ouv}{\(\mathrm{O};\vec{u},\vec{v}\)}
\providecommand{\intervFF}[2]{\ensuremath{[#1;#2]}}
\renewcommand{\euro}{\ifmmode\text{\texteuro}\else\texteuro\fi}
\DeclareUnicodeCharacter{00B0}{\ifmmode^{\circ}\else\textdegree\fi}
\DeclareUnicodeCharacter{20AC}{\ifmmode\text{\texteuro}\else\texteuro\fi}
% Énoncés : Éducation nationale. Transcription et corrigés : APMEP, Denis Vergès et contributeurs. Découpage et classement : CoopMaths, Sébastien Lozano, Éric Elter et contributeurs. Adaptation de la présentation : Maths Explorer.
% CC BY-SA : https://coopmaths.fr/www/about/
% https://coopmaths.fr/alea/static/eam/2026/tex/eam_specialite_2026_09_metropole_1_cor.tex
\begin{document}
\noindent\textit{APMEP / CoopMaths. CC BY-SA. Présentation adaptée par Maths Explorer.}\par\medskip
\emph{On note $E$ l'évènement \og le patient est un enfant \fg{} et $M$ l'évènement \og le médecin prescrit des
médicaments \fg. On a $P(E) = 0,7$, $P_E(M) = 0,2$ et $P_{\overline{E}}(M) = 0,4$.}

\begin{enumerate}
	\item La somme des probabilités issues d'un
	même n{\oe}ud vaut 1 : $P(\overline{E}) = 1 - 0,7 = 0,3$,
	$P_E(\overline{M}) = 1 - 0,2 = 0,8$ et $P_{\overline{E}}(\overline{M}) = 1 - 0,4 = 0,6$.

	\begin{center}
		\begin{tikzpicture}[x=2.8cm, y=0.9cm]
			
			\node (root) at (0,0) {};
			\node (E) at (1,1.2) {$E$};
			\node (Eb) at (1,-1.2) {$\overline{E}$};
			\node (M1) at (2,1.9) {$M$};
			\node (M1b) at (2,0.5) {$\overline{M}$};
			\node (M2) at (2,-0.5) {$M$};
			\node (M2b) at (2,-1.9) {$\overline{M}$};
			
			\draw (root.east) -- (E) node[pos=0.5,above] {0,7};
			\draw (root.east) -- (Eb) node[pos=0.5,below] {0,3};
			
			\draw (E.east) -- (M1) node[pos=0.5,above] {0,2};
			\draw (E.east) -- (M1b) node[pos=0.5,below] {0,8};
			\draw (Eb.east) -- (M2) node[pos=0.5,above] {0,4};
			\draw (Eb.east) -- (M2b) node[pos=0.5,below] {0,6};
		\end{tikzpicture}
	\end{center}

	\item La probabilité d'un chemin de l'arbre est le produit des probabilités rencontrées le
	long de ce chemin :
	\[P(\overline{E} \cap M) = P(\overline{E}) \times P_{\overline{E}}(M) = 0,3 \times 0,4 = 0,12.\]
	La probabilité que le patient choisi soit un \textbf{adulte} \textbf{et} que le médecin lui prescrive des médicaments est de $0,12$ ;
	autrement dit, 12\,\% des patients sont des adultes à qui le médecin prescrit des médicaments.
	

	\item $E$ et $\overline{E}$ forment une partition de l'univers. D'après la formule des
	probabilités totales :                          
	
	 \[ P(M) = P(E \cap M) + P(\overline{E} \cap M) \]
 \[ P(M) = P(E) \times P_E(M) + P(\overline{E} \cap M) \]
 \[ P(M) = 0,7 \times 0,2 + 0,12 \]
 \[P(M) =  0,26\]
	
	La probabilité que le médecin prescrive des médicaments au patient est égale à $0,26$.
	

	\item Deux évènements $E$ et $M$ (avec $P(E) \neq 0$) sont indépendants si et seulement si
	
	$P(E \cap M) = P(E) \times P(M)$.

	D'une part $P(E \cap M) = 0,14$ ; d'autre part
	$P(E) \times P(M) = 0,7 \times 0,26 = 0,182$.

	Comme $0,14 \neq 0,182$, l'égalité n'est pas vérifiée. 

	Les évènements $E$ et $M$ ne sont \textbf{pas indépendants}.


	\item On cherche la probabilité conditionnelle $P_M(E)$.
	\[P_M(E) = \dfrac{P(E \cap M)}{P(M)} = \dfrac{0,14}{0,26} = \dfrac{14}{26} = \dfrac{7}{13}.\]
	Sachant que le médecin a prescrit des médicaments, la probabilité que le patient soit un
	enfant est égale à $\dfrac{7}{13}$.
\end{enumerate}

\end{document}
